\section{Results}
\label{sec:results}

Every number in this section is produced by a command recorded in the replication package's run log
and, in the manuscript, by a generated table of macros built from the released result files. No value is typed by hand.

\subsection{The frame}

The enumeration produced \Nframe{} unique repositories from \Ntopicpages{} topic pages spanning
\topicslugs{} topic slugs and \langfilters{} language filters, with the SHA-256 of the concatenated
page payload recorded in the released frame manifest. The frame is a snapshot: re-running the
enumeration later would produce a different membership, which is why the snapshot timestamp and hash
are part of the replication package.

\subsection{Static screening and attrition}

Table~\ref{tab:funnel} reports the screening funnel. Of \Nscreened{} repositories screened,
\Neligible{} satisfied all eight eligibility criteria, a yield of \YieldPct\,\% (95\,\% Wilson interval
\YieldLo--\YieldHi\,\%).

\begin{table*}[t]
\centering
\caption{Static screening funnel. Every screened repository is counted exactly once. Reasons are
ordered so that the recorded reason is the most fundamental blocker.}
\label{tab:funnel}
\begin{tblr}{width=\textwidth,colspec={X[2.2,l]rr},rowsep=1pt}
\toprule
\textbf{Stage / reason} & \textbf{n} & \textbf{share} \\
\midrule
Repositories screened            & \Nscreened  & 100.0\,\% \\
\quad Excluded: too few scenarios & \RtooFew   & \PtooFew\,\% \\
\quad Excluded: no feature files  & \RnoFeat   & \PnoFeat\,\% \\
\quad Excluded: no step definitions & \RnoSteps & \PnoSteps\,\% \\
\quad Excluded: no BDD runner dependency & \RnoRunner & \PnoRunner\,\% \\
\quad Excluded: no application shape & \RnoShape & \PnoShape\,\% \\
\quad Excluded: licence not permissive & \Rlicence & \Plicence\,\% \\
\quad Excluded: language out of scope & \Rlang & \Plang\,\% \\
\quad Excluded: no build config & \RnoBuild & \PnoBuild\,\% \\
\quad Excluded: other / fetch failure & \Rother & \Pother\,\% \\
\midrule
\textbf{Eligible}                 & \Neligible & \YieldPct\,\% \\
\bottomrule
\end{tblr}
\end{table*}

The dominant exclusions are substantive rather than technical. ``Too few scenarios'' is the single
largest category, which indicates that most topic-tagged repositories are small demonstrations or
single-scenario examples. ``No feature files'' is the second largest: a substantial share of
repositories carrying a BDD topic tag do not contain a single \texttt{.feature} file, which means
topic tags are an unreliable indicator of BDD use even at the level of file presence. Together these
two categories account for the great majority of exclusions.

The application-shape criterion, added after the pilot, excluded a further \RnoShape{}
repositories that had features, steps, a build and a licence but fewer than ten non-test source files.
These are the boilerplates, formatters, step libraries and reporting agents that pass a naive
eligibility test while having no business-rule domain to mutate. Their exclusion is the clearest
single demonstration of why the naive test is inadequate for this construct.

\subsection{Build and execution screen}

Static eligibility is necessary but not sufficient: a repository must also build and actually execute
its Gherkin. Table~\ref{tab:build} reports the outcome of the build and execution screen over the
\Nbuildtested{} eligible repositories that were attempted.

\begin{table*}[t]
\centering
\caption{Build and execution screen. \texttt{no\_bdd\_execution} denotes a repository whose tests pass
while nothing indicates that the Gherkin layer ran; it is not counted as a build failure, and it is
critically not counted as a pass.}
\label{tab:build}
\begin{tblr}{width=\textwidth,colspec={X[2.4,l]rX[1.6,l]},rowsep=1pt}
\toprule
\textbf{Outcome} & \textbf{n} & \textbf{Representative reasons} \\
\midrule
\texttt{bdd\_passes} (green Gherkin baseline) & \Ngreen & --- \\
\texttt{bdd\_fails} (tests ran, failed) & \Nbdffails & failing scenarios at HEAD \\
\texttt{no\_bdd\_execution} & \Nnobddexec & tests pass, no Cucumber evidence \\
\texttt{timeout} (exceeded the run budget) & \Ntimeout & repository too slow to screen within budget \\
\texttt{unsupported} (no runner available) & \Nunsupported & language/tool chain out of scope \\
\texttt{build\_fails} & \Nbuildfails & \texttt{dependency\_rot},
\texttt{needs\_browser}, \texttt{compile\_error}, \texttt{tool\_incompatibility} \\
\midrule
\textbf{Attempted} & \Nbuildtested & of \Neligible{} eligible \\
\bottomrule
\end{tblr}
\end{table*}

Three observations follow directly from this table and are important for the interpretation of the
whole study.

First, \textbf{a green \texttt{mvn test} or \texttt{npm test} does not prove that a BDD suite ran}. The
\texttt{no\_bdd\_execution} category exists because we initially accepted a zero exit code as evidence
of BDD execution and a spot check showed this to be false: a repository whose Maven build succeeded
had executed none of its Gherkin. The harness was changed to require positive evidence of Cucumber
execution in the output before accepting a repository, and \Nnobddexec{} attempted repositories fell
into this category. Any study that equates ``tests pass'' with ``BDD suite passes'' will overcount its
corpus.

Second, \textbf{the failures are heterogeneous}, and conflating them would misstate the evidence.
\texttt{dependency\_rot} is a property of the corpus (old pins no longer install);
\texttt{needs\_browser} is a property of the harness environment, not of the suite's quality; a
failing scenario at HEAD is a property of the repository. Only the last is a statement about the
suite's health, and only the first two are avoidable by a better-equipped environment.

Third, \textbf{the rate of green, executing BDD baselines in this population is low}. We report the
count rather than a rate, because the build screen did not complete for every eligible repository: of
\Neligible{} eligible repositories, \Nbuildtested{} were attempted. The remaining repositories were
not screened for build and execution within the resources of this study, and we therefore do not
report a green-suite \emph{rate} that would imply full coverage of the eligible set. This is a
limitation of our execution budget, not a finding.

\subsection{Tool feasibility and instrument validation}

Before any corpus measurement we established, empirically, that mutation analysis can be driven by the
BDD runner itself and that the three-mode decomposition of
Section~\ref{sec:method} is computable. Table~\ref{tab:feasibility} reports the harness
demonstrations. Each suite now contains a deliberately unexercised domain module, so the
\texttt{NO\_COVERAGE} column is a measured control rather than an assumption.

\begin{table*}[t]
\centering
\caption{Three-mode mutation measurement on the purpose-built feasibility suites. Both suites
contain a domain module that no test references, which is what makes the \texttt{NO\_COVERAGE} column
a real control: a non-zero count demonstrates that the tool distinguishes \emph{executed and
survived} from \emph{never executed}, and that $\msall$ and $\mscov$ are therefore separate
quantities. $S$ is the comparison set of Section~\ref{sec:method}. The Java figures are exact because
PIT reports coverage in every mode; the JavaScript \dk{} is an \emph{upper bound}, because Stryker's
command runner cannot report per-test coverage for the non-BDD mode, so $S$ collapses to the
BDD-covered set and the non-BDD kill set is a lower bound. Toolchains are pinned: Stryker \StrykerVer{} with cucumber-runner \CukeRunnerVer{} and Cucumber \CucumberJsVer{} for JS/TS; PIT \PitVer{} with pitest-junit5-plugin \PitPluginVer{} and Cucumber-JVM \CucumberJvmVer{} for Java. Operator alignment was applied to PIT only: the JavaScript harness uses Stryker's default operators throughout, so its two mutant-count rows coincide by construction rather than by coincidence.}
\label{tab:feasibility}
\begin{tblr}{width=\textwidth,colspec={X[2.6,l]X[0.9,r]X[0.9,r]},rowsep=1pt}
\toprule
\textbf{Quantity} & \textbf{JS/TS} & \textbf{Java} \\
\midrule
Mutants, aligned decision-logic operators & \JsMutants & \JavaMutants \\
Mutants, tool defaults & \JsMutants & \JavaDefaultMutants \\
BDD-only: killed & \JsKilled & \JavaKilled \\
BDD-only: \texttt{NO\_COVERAGE} & \JsNoCov & \JavaNoCov \\
BDD-only $\msall$ (\%) & \JsMsAll & \JavaMsAll \\
BDD-only $\mscov$ (\%) & \JsMsCov & \JavaMsCov \\
\midrule
Comparison set $|S|$ & \JsS & \JavaS \\
$\dkbdd$ (BDD unique uplift) & \JsDkBdd{} (upper bound) & \JavaDkBdd{} (exact) \\
$\dknon$ (non-BDD unique uplift) & \JsDkNon & \JavaDkNon \\
Overlap & \JsOverlap & \JavaOverlap \\
$\bddunique$ & \JsBddUnique & \JavaBddUnique \\
$\bddtotal$ & \JsBddTotal & \JavaBddTotal \\
\bottomrule
\end{tblr}
\end{table*}

\textbf{The \texttt{NO\_COVERAGE} control.} The BDD-only runs leave \JsNoCov{} of \JsMutants{}
JavaScript mutants and \JavaNoCov{} of \JavaMutants{} Java mutants unexecuted. Both counts are
non-zero, so the tools demonstrably separate execution from non-execution, and
$\msall$ (\JsMsAll\,\% and \JavaMsAll\,\%) differs from $\mscov$ (\JsMsCov\,\% and
\JavaMsCov\,\%) on the same runs by more than a factor of two. A run reporting a single pooled score
would therefore have understated the discriminating power of these suites by half.

\textbf{What the decomposition reveals, and why the correction mattered.} On the Java suite the
specification layer kills \JavaKilled{} of the \JavaS{} mutants in $S$, a healthy $\mscov$ of
\JavaMsCov\,\%. Its unique uplift is nevertheless $\dkbdd = \JavaDkBdd$: every mutant it kills, the
unit tests also kill. The two mutants that only the unit tests kill are boundary cases the scenarios
never probe, giving $\dknon = \JavaDkNon$. Under the earlier, inverted definition
$\mscov(\full,\mathcal{S})-\mscov(\bdd,\mathcal{S})$ this repository would have been reported as
having a BDD uplift of $+\JavaDkNon$ --- the exact opposite conclusion, since the number is produced
entirely by the unit tests. This is a small, fully known suite, so the result is illustrative rather
than general; but it is precisely the configuration the metric has to get right, and the inverted
version got it wrong on the first suite we tried it on.

The JavaScript suite behaves differently: $\dkbdd = \JsDkBdd$ and $\bddunique = \JsBddUnique$, so a
small number of mutants are killed only by the scenarios. We do not over-read this, because the
figure is an upper bound as explained in the caption.

\textbf{Operator alignment.} Aligning both tools to a decision-logic operator set changes the Java
mutant count from \JavaDefaultMutants{} (tool defaults) to \JavaMutants{} (aligned), a difference of
one mutant on this suite. The cross-tool contrast is much larger: PIT generates \JavaMutants{}
mutants where Stryker generates \JsMutants{} for comparably sized domain logic, a ratio of roughly
\MutantRatio{}. That is consistent with the cross-tool incomparability reported by
\citet{Kintis2018HowEffectiveTools}, and it is why the analysis refuses to pool scores across
ecosystems. Aligning operators makes the comparison \emph{controlled}; it does not make the two
scales interchangeable.

\subsection{Known-ground-truth validation: seeded drift}\label{sec:seeded}

Feasibility is not sensitivity. To test whether the seven indicators respond to drift that is
actually present, we injected one controlled drift mutation at a time into the working JavaScript
feasibility suite, re-measured every indicator and the BDD-only mutation score, and reverted the tree
between injections. The baseline BDD-only $\mscov$ is \DriftBaseMsCov\,\%. An injection was accepted
only if the unmutated suite still passed afterwards, so a broken experiment can never be mistaken for
an insensitive indicator; \DriftInvalid{} of the six injections failed that check and is reported
rather than analysed.

Of the \DriftTargets{} sensitivity targets, \DriftDetected{} behaved as designed. Tagging scenarios as
skipped raised $\mathit{D1}$ by \DriftIoneDelta{} with \emph{no} change in mutation score, which is
correct: the skipped scenarios were killing nothing. Duplicating a scenario raised
$\mathit{D4}$ by \DriftIthreeDelta{}, again with no change in the score, since a duplicate detects
nothing new. The negative control, which only reformatted the suite and added comments, moved nothing
at all --- the test that an indicator cannot pass merely by always answering ``drift''.

The two misses are the interesting ones, and they are both about the oracle rather than the text.
Replacing an assertion with a comparison of the computed value against itself, so that it can never
fail, left $\mathit{D2}$ \emph{unchanged} while the mutation score collapsed from
\DriftBaseMsCov\,\% to \DriftItwoMsCov\,\%. Rewriting the expected value so that it is re-derived from
the production calculation --- a mirror step, the exact D7 construct --- left $\mathit{D7}$
\emph{unchanged} while the score fell to \DriftIfourMsCov\,\%. In both cases the indicator that is
supposed to detect the flaw did not move, and the mutation score lost more than half its value.

We draw the honest conclusion rather than the flattering one. The static indicators are not
interchangeable substitutes for mutation-based measurement, and this study does not use them as such:
$\mathit{D2}$ and $\mathit{D7}$ enter the analysis as covariates and early-warning signals, and the
mutation-derived quantities remain the outcome. The campaign also bounds what can be claimed for the
indicator validity in general: two of four targets detected on one suite is a demonstration that the
instrument can detect seeded drift, not an estimate of its sensitivity, and a sensitivity estimate
would need many more injections across several suites. Finally, the wrong-literal injection shows a
ceiling the indicators cannot cross: $\mathit{D6}$ counts whether a literal is present, not whether it
is \emph{correct}, so a suite with a confidently wrong literal looks identical to one with a right one
until the test is executed.

\subsection{The eligible population, described}\label{sec:census}

RQ1 also carries a descriptive obligation that requires no builds at all: the distribution of the seven
indicators across every statically eligible repository. Omitting it, as the first version of this paper
did, left the measurement half-evaluated even though the data were free. We measured
\CensusMeasured{} of the \CensusEligible{} eligible repositories; the remaining seven are large
monorepos that exceeded a fixed per-repository time budget and are recorded as timeouts rather than
silently dropped. Table~\ref{tab:census} reports the distributions.

\begin{table}[!htbp]
\centering
\footnotesize
\caption{Static indicator census over the statically eligible population. $n$ is the number of
repositories for which the indicator is defined; a repository where an indicator is undefined is
excluded from that row and counted, never imputed as zero. D3 is undefined throughout because it
requires version-control history and the working copies carried none.}
\label{tab:census}
\begin{tblr}{width=\columnwidth,colspec={lrrr},rowsep=1pt}
\toprule
\textbf{Indicator} & \textbf{$n$} & \textbf{median} & \textbf{IQR} \\
\midrule
D1 skip ratio & \CensusDoneN & \CensusDoneMedian & [\CensusDoneQone, \CensusDoneQthree] \\
D2 assertion weakness & \CensusDtwoN & \CensusDtwoMedian & [\CensusDtwoQone, \CensusDtwoQthree] \\
D3 stale specification & \CensusDthreeN & \CensusDthreeMedian & --- \\
D4 redundancy clusters & \CensusDfourN & \CensusDfourMedian & [\CensusDfourQone, \CensusDfourQthree] \\
D5 glue duplication & \CensusDfiveN & \CensusDfiveMedian & [\CensusDfiveQone, \CensusDfiveQthree] \\
D6 literal-bearing steps & \CensusDsixN & \CensusDsixMedian & [\CensusDsixQone, \CensusDsixQthree] \\
D7 mirror steps & \CensusDsevenN & \CensusDsevenMedian & [\CensusDsevenQone, \CensusDsevenQthree] \\
\bottomrule
\end{tblr}
\end{table}

Three features of Table~\ref{tab:census} matter for the argument. First, D1 has essentially no
variance: the median suite marks almost no scenario skipped (median \CensusDoneMedian, upper quartile
\CensusDoneQthree), so as a discriminator in \emph{this} population the indicator is close to useless,
whatever its diagnostic value in a team where skipping is common. Second, D2 spreads widely --- median
\CensusDtwoMedian{} with an upper quartile of \CensusDtwoQthree{} --- so a quarter of these suites have
more than half of their step definitions classified as weak or non-asserting by the heuristic.
Third, D6 is saturated at \CensusDsixMedian: essentially every step definition carries a literal, which
is consistent with the ceiling the seeded-drift campaign found for D6 independently.

Two readings must be resisted. A D7 median of \CensusDsevenMedian{} does \emph{not} establish that
mirror steps are absent from public suites: the seeded-drift campaign showed that D7 failed to detect a
mirror step built by hand, so a zero here is weak evidence of absence rather than evidence of a weak
threat. And none of these distributions is an association: \CensusDtwoN{} repositories have D2 defined,
and \CensusLowConf{} of them are flagged low-confidence by the extractor itself, so the census
describes the population and nothing more. No correlation with oracle strength is estimated, because
no repository reached the mutation stage.

\subsection{Oracle strength (RQ2)}

RQ2 asks what share of covered domain mutants the BDD suite alone kills, and how much it adds beyond
the repository's other tests. The instrument to answer this is demonstrated above. On the corpus, the
answer is bounded by what could be measured: \Nmeasured{} repositories yielded a complete
BDD-only measurement, and \Nwithfull{} of those also yielded a full-suite comparison from which
$\dk$ could be computed on a common mutant set.

\MeasuredTable

We report this plainly rather than substituting a number that the data do not support. The construct
is operational and the harness computes it; the population did not supply the units of analysis at the
scale the design anticipated. Section~\ref{sec:discussion} draws out why, and
Section~\ref{sec:threats} states what a reader should and should not infer.

\subsection{Association and erosion (RQ3)}

RQ3 asks whether pre-specified static drift indicators associate with lower BDD-only kill rates, and
whether drift accumulates over repository history.

The protocol contains its own rule for this situation. The protocol's drop rule DR1 removes a stratum
with fewer than eight repositories measurable end-to-end. Here \Nmeasured{} repositories were
measurable, so \textbf{both strata fall below the threshold and DR1 applies to each}: the confirmatory
family (D1 and D2 against $\mscov(\bdd)$, Holm-corrected) is therefore \textbf{void by the protocol's own
rule}, not merely underpowered. We invoke DR1 rather than reporting a null result, because a null
would imply that an association had been estimated and found absent, and no estimation was possible.
The pre-specified detectable-effect floor is consequently not reported: it is defined on the Fisher-$z$
scale as $1.030/\sqrt{n-3}$, which is undefined for the $n$ actually achieved.

For the erosion component, the static history miner recomputes the indicators at \HistoryK{}
deterministically sampled commits per repository and reports Theil--Sen slopes. This component needs
no builds, but it does need commit history, and the repositories that reached this stage were
materialised from source archives and from depth-one clones, neither of which carries the history the
miner walks. The series is therefore empty for a \emph{mechanical} reason, not because the
repositories lack history: \Nhistory{} repositories yielded a usable series. The miner is implemented
and shipped in the replication package, and exercising it requires re-cloning candidates with full
history --- a step this study budgeted for the corpus stage and did not reach. We make no longitudinal
claim.

\HistoryTable

\subsection{What could not be measured, and why it matters}

We close the results with an explicit inventory of what this study did not establish, because a reader
is entitled to know where the evidence stops.

\begin{itemize}
\item \textbf{No inter-rater reliability statistic.} The two judgement-coded indicators (D2's weakening
      rubric and D7's mirror judgement) require independent human raters. None were available. No
      $\kappa$ is reported, and no automated labeller was substituted, because agreement between a
      heuristic and itself is not inter-rater reliability.
\item \textbf{No validated precision for D7.} D7 is reported only as an automated census of its loose
      conjunction.
\item \textbf{No completed build screen.} \Nbuildtested{} of \Neligible{} eligible repositories were
      attempted; the remainder were not screened for build and execution.
\item \textbf{No corpus-scale association.} The pre-specified confirmatory family is void under drop
      rule DR1, because \Nmeasured{} repositories were measurable against a threshold of eight per
      stratum. No association was estimated, so no null result is claimed either.
\item \textbf{No Python stratum.} The frame yielded \NpyEligible{} Python repositories satisfying the
      application-shape criterion, and the Python mutation toolchain could not be established
      reliably in this environment (package index unreliability, and the tooling's coupling to a
      unit-test runner rather than a BDD runner). The planned three-stratum design was therefore not
      realised, and language effects could not be separated from tool effects.
\end{itemize}

